% English translation of questions/5780/semB-moedB/q2a.tex (metadata is taken from the Hebrew file).

\begin{question}
(15 pts) For which values of the parameter $a$ does the equation $(x-1)^2e^x-a=0$ have exactly one real solution?
\end{question}

\begin{hint}
Investigate the function $g(x)=(x-1)^2e^x$ (intervals of monotonicity, extrema, limits at $\pm\infty$), and count how many times the horizontal line $y=a$ intersects the graph.
\end{hint}

\begin{solution}
The equation is equivalent to $g(x)=a$ for $g(x)=(x-1)^2e^x$. We investigate $g$, which is differentiable on all of $\R$:
\[ g'(x)=2(x-1)e^x+(x-1)^2e^x=e^x(x-1)(x+1). \]
$e^x>0$, hence $g'>0$ on $(-\infty,-1)$, $g'<0$ on $(-1,1)$ and $g'>0$ on $(1,\infty)$. That is, $g$ is strictly increasing on $(-\infty,-1]$, strictly decreasing on $[-1,1]$ and strictly increasing on $[1,\infty)$, with a local maximum $g(-1)=\frac{4}{e}$ and a local minimum $g(1)=0$.

Limits: $\lim_{x\to\infty}g(x)=\infty$; and as $x\to-\infty$, $(x-1)^2e^x=\frac{(x-1)^2}{e^{-x}}\to0$ (the exponential dominates a polynomial, e.g. by L'Hôpital twice). Also $g(x)\ge0$ for every $x$, with equality only at $x=1$.

By the Intermediate Value Theorem and strict monotonicity on each interval, on each interval of monotonicity every value in the range of that interval is attained exactly once. The ranges of the intervals: $(-\infty,-1]\mapsto(0,\frac4e]$, $[-1,1]\mapsto[0,\frac4e]$, $[1,\infty)\mapsto[0,\infty)$. Hence the number of solutions:
\begin{itemize}
  \item $a<0$: no solutions ($g\ge0$).
  \item $a=0$: a unique solution, $x=1$.
  \item $0<a<\frac4e$: three solutions (one in each interval of monotonicity).
  \item $a=\frac4e$: two solutions ($x=-1$, and one more in $(1,\infty)$).
  \item $a>\frac4e$: a unique solution (in $(1,\infty)$; on the other two intervals $g\le\frac4e<a$).
\end{itemize}
\textbf{Answer:} the equation has a unique real solution for $\boxed{a=0 \text{ or } a>\frac{4}{e}}$.
\end{solution}
